\section{Adaptive Pool Configuration}
请求调度处理单次到达，自适应配置则决定后续一段时间的服务能力。控制器根据近期到达率、输入输出长度及尚未完成的工作，联合选择 P、D、M 实例数、各角色频率、停车状态和 PD/M 分流阈值。以 $c$ 表示配置，$H$ 表示计划保持窗口，其目标为
\begin{equation}
 c^*=\arg\min_{c\ \mathrm{feasible}}
       \{H P(c)+E(c_0,c)\},
 \label{eq:pool-objective}
\end{equation}
其中 $P(c)$ 包含工作与停车实例的功率，$E(c_0,c)$ 是从当前配置转换的额外能量。可行候选必须满足带余量的 TTFT、TPOT 约束、KV 容量和 profiler 覆盖范围，且保留仍有在途请求的服务分支。规划器对每个候选分别计入转换成本；当前配置可行时，只有收益超过相对阈值才切换。默认窗口为 $60$ 秒、阈值为 $3\%$，实际可由运行配置覆盖。转换成本优先采用绑定配置与 profile 的配对测量；严格模式排除缺少证据的转换，普通模式保留显式标记的估计回退。

负载输入不仅包含新请求速率，也包含等待 prefill 的 token 数、剩余 decode 工作量和已经占用的 KV。长请求即使离开近期统计窗口，也继续计入在途快照。控制器默认每秒观察保护信号、每十秒进行普通规划；测量覆盖不足时，维持能够解释的当前配置或进入保守恢复路径，而非把无法预测的候选当成节能机会。

\subsection{Frequency Control}
频率控制在保持模型与 TP 拓扑不变的条件下改变服务速率。规划器枚举 profiler 支持的离散频点，根据各阶段的负载、预测时延和功率选择 $f_P,f_D,f_M$。P 与 D 可以使用不同频率，从而分别适应 prefill 与 decode 的资源需求；同一实例的全部 TP GPU 使用一致的目标频率。当前搜索是离散候选枚举，实际执行只修改状态发生变化的实例。

周期规划之外，启用预算感知保护器时，快速路径根据压力来源定位应调整的角色：PD 的 prefill 压力作用于 P，decode 压力作用于 D，混合路径压力作用于 M。其首级动作提高相关角色频率，持续压力再触发容量恢复；短暂的集体 token 间隙只允许一次有界调频探测。启用截止时间保护时，控制器还会在较高频点中寻找能够消除当前风险的最低频率，并为 M 维持频率下限，直到连续安全观察后释放。降频同时受保持时间、冷却窗口和多次预测确认约束，避免噪声引起反复切换。

物理执行以 GPU 标识加互斥锁，防止不同控制动作同时写入同一张卡；阻塞的硬件调用放入工作线程，使正在加载的实例不会阻塞其他实例的紧急调频。这里的共享资源约束落在物理 GPU，而 profile 的频率域表示获得测量支持的频点集合，二者分别核查。

\subsection{Role Change}
角色控制改变各阶段获得的实例数量。规划器在满足总实例预算和服务约束的布局中选择 P/D/M 比例，并可将剩余实例置为驻留低功耗状态或停止进程。将目标数量映射到具体实例时，首先保留已有角色，再结合驻留状态与负载安排变更，以减少不必要的唤醒和角色重分配。

当前实现区分路由角色调整与释放资源。P、D、M 活跃角色之间的转换仅更新后续请求的路由角色及频率；已进入系统的请求继续绑定原实例，不迁移在途 KV。因此，角色改变可以渐进生效，过渡期间实例仍需完成先前接收的工作。规划时的 backlog 保留其原有分支归属，不能因为目标角色改变而删除存量负载。

当实例进入停车或停止状态时，控制器先关闭新请求准入，再等待代理侧 prefill 与 decode 工作清零，并在接入原生控制时核查全部 rank 的队列、KV 块及待完成传输均已释放，然后执行降功耗或停止。唤醒则完成启动或恢复、就绪确认、调频和原生准入恢复后才向路由器发布。多个实例的动作可以并行，但任何失败都阻止受影响实例继续接收请求。由此，轻量角色变化保留服务连续性，资源回收具有明确的排空边界。

M 实例数量还受经验容量下限约束。降低该下限需要与模型、TP、频率、SLO 和工作负载域匹配的验收记录；负载退出已验证范围时，控制器优先恢复预留 M 容量。该约束补足排队近似对复杂混合干扰的描述，避免仅凭理想化预测过度缩容。

\subsection{TP Weight Change}
当前 TP 权重改变指两个\emph{常驻异构 TP 池之间的请求份额}。两个池使用不同的 TP 大小和不重叠的 GPU，各自绑定独立 profile；在线调整保持既有模型分片布局。外层规划枚举两池权重 $w$ 与 $1-w$，内层分别求解角色和频率配置，并按式~\eqref{eq:pool-objective} 汇总两池能量。默认权重步长为 $0.1$，候选还包括当前目标份额和上一周期的实际份额。

对于每个权重，未来到达率按份额拆分，而已进入系统的请求仍计入原池 backlog。即使某池权重降为零，其全部驻留副本和空闲 GPU 仍参与功耗核算，防止将备用容量视为免费资源。仅当两池均满足覆盖、容量和 SLO，且新的联合目标达到节省阈值时，协调器才执行各池配置；全部角色、频率和准入状态就绪后发布新权重。

联合规划还要求可用的实测 decode 与混合窗口功率。保护器处于扩容或保持阶段时暂缓新的能量优化；若任何池执行失败，保持原目标权重并暂停后续优化，等待真实状态恢复。这样，流量份额发布始终发生在实际容量调整之后，避免请求先被引向尚未就绪的池。

在线分流采用份额欠额：若已分配 $N$ 个请求、池 $j$ 获得 $n_j$ 个，则在当前可行池中选择
\begin{equation}
 j^*=\arg\max_j\{w_j(N+1)-n_j\}.
\end{equation}
目标池不可行时可使用另一可行池，因此份额是控制目标而非强制拒绝条件。实际分配数反馈到下一周期，使联合规划能观察容量约束导致的偏离。此机制不要求跨 TP 搬移在途请求；在线模型权重重分片入口目前仍因缺少经过 GPU 验证的原生事务后端而关闭。
